\scanpage{5}
\uline{\unsure{Rere|Re re|Rem} NB} On aura \unsure{néc.|même}$^{\mathrm t}$ $h_0\in\Gamma(C\smallsetminus Q)$ ....

\margin{non autodual sauf si $\beta$ inv.}
$C^\infty_{vect\,orth}$ cat. des syst. $(V,f,\varphi_G^V,a_0)$, où $V$ fibré vectoriel de rg $n+1$, $f$ forme quadratique \deleted{non id. nulle sur aucune fibre}, $\varphi_G^{V\,orth}:\unsure{\mathscr G|G|\mathfrak G}\longrightarrow\operatorname{Aut}(V,f)$, $a_0$ section de $V$, satisfais. à

a) les $\tau_i^V$ réflexions orth \inserted{$(0\leq i\leq n)$}, les $\tau_i$ $(1\leq i\leq n)$ \deleted{\unsure{np|non}} $\ne$ id en chaque fibre ($\Rightarrow L_i,H_i$ pour $1\leq i\leq n$), $\tau_0$ associé à $\mathcal O_S a_0$, $\deleted{(a_0\ne0\text{ en tte}}$ et $\boxed{f(a_0)=1}$ (de sorte que $L_0$ $\mathcal O_S\deleted{\unsure{b_0}}a_0$ ss-fibré vect\ldots)

b), c) comme dessus, avec commentaire correspondant.

On pourrait définir $\widetilde C^\infty_{vect\,orth}$, catégorie des \deleted{systèmes} $(V,f,\varphi_G^V,L_0)$, où \deleted{$(\unsure{V},\unsure{f},\varphi_G^V,L_0)$} \deleted{\unsure{comme} $\varphi_G^V$, $f$} $(V,f)$ fibré vectoriel quadratique \deleted{comme dans $C^\infty_{vect}$, $(V,f)$ fibré vect.} de rg $n+1$, $\varphi_G^V:\mathscr G_n\longrightarrow\operatorname{Aut}(V,f)$, $L_0$ ss-fibré de rg $1$, tels que a) b) c) comme précédemt, mais $f(a_0)=1$ remplacé par $f\mid L_0$ est \uline{\unsure{lisse|libre}}.

Et On trouve $C^\infty_{vect\,orth}$ à partir de $\widetilde C^\infty_{vect\,orth}$ par rigidification par une section $a_0$ de $L_0$ telle que $f(a_0)=1$ (les sections forment un torseur sous $\mu_2$ \ldots\ ).

\diagram{\begin{tikzcd}[ampersand replacement=\&,column sep=small]
C_{pr}\arrow[r,dash,"\approx"]\& C'_{pr}\arrow[r,dash,"\approx"]\& C_{aff}\arrow[r,dash,"\approx"]\& C'_{aff}\arrow[r,dash,"\approx"]\& C_{vect}\arrow[r,dash,"\approx"]\& C'_{vect}\,\Big|\arrow[r,dash,"\approx"]\& C^\infty_{pr}
\end{tikzcd}}
\margin{$n\geqslant2$ (expliciter variante si $n=1$ !)}
\diagram{\begin{tikzcd}[ampersand replacement=\&,column sep=small]
{}\arrow[r,dash,"\approx"]\& C^\infty_{vect}\arrow[r,dash,"\approx"]\& \Big|\,C^\infty_{pr\,orth}\arrow[r,dash,"\approx"]\& C^\infty_{vect\,orth}
\end{tikzcd}}
\margin{si de plus les $\lambda_i=-1$ i.e. les $\tau_i$ des réflexions.}

% archive sheet number: 4

